\chapter{数学与理论}
\label{ch:MathAndTheory}

本章介绍描述电磁波以及理解式~\ref{eq:GNLSE} 所需的基本数学工具。


\section{实场与复场}
%Explain that the real field is what determines how an electron moves and that the complex one is only used for mathematical convenience when dealing with phase shifts.
一个电荷为 $q$、质量为 $m$ 的粒子若处于电场 $\Bar{\E}_r=\E_r \hat{x}$ 中，其加速度为
\begin{align}
    \Bar{a} &=  \frac{q\E_r}{m}\hat{x}.
\end{align}
电场上的下标 $r$ 表明这是真实的“实电场”，它决定带电粒子究竟如何加速。相比之下，“复电场”只是一个数学工具，它能让涉及电磁波的计算更容易处理，而实电场总能从复电场中恢复出来。举例来说，在体介质中传播的实电场波，如果其空间角频率 $\betag$ 依赖于时间角频率 $\omega$，则可以写成
\begin{align}
\label{eq:real_field}
    \E_r(z,t) &= |\E_0|\cos\left(\betag(\omega)z-\omega t+\phi\right) \\\nonumber
     &=|\E_0| \real\left\{  \exp\left( i\betag(\omega)z-i\omega t+i\phi \right) \right\} \\ \nonumber
     &= \real\left\{ \E_0 \exp\left( i\betag(\omega)z-i\omega t\right) \right\}  \\ \nonumber
     &= \real\left\{ \E(z,t) \right\}  \\ \nonumber
     &=  \frac{1}{2}\left( \E(z,t) + \E^*(z,t) \right).
\end{align}
关于式~\ref{eq:real_field} 的图像化说明，可参考 \href{https://www.desmos.com/calculator/fgvozursrl}{这个交互图}。

在建模相位变化和干涉时，使用复电场尤其方便。比如，要给实电场引入一个相移 $\phi_0$，就必须“手动”把它写进余弦函数：
\begin{align}
\label{eq:insert_phase}
    \E_r(z,t) &= |\E_0|\cos\left(\betag(\omega)z-\omega t+\phi\right) \Rightarrow \\ \nonumber \E_r'(z,t) &= |\E_0|\cos\left(\betag(\omega)z-\omega t+\phi+\phi_0^{\textcolor{Red}{\swarrow}}\right).
\end{align}
而在复场表示中，同样的操作只需做一次复数乘法，因为
\begin{align}
\label{eq:insert_phase_complex}
    \E(z,t) &=\E_0 \exp\left( i\betag(\omega)z-i\omega t+ i\phi\right)\Rightarrow \\ \nonumber \E'(z,t) &=  \E_0 \exp\left( i\betag(\omega)z-i\omega t +i\phi\right)\exp\left( i\phi_0\right) \\ \nonumber
    &=\E_0 \exp\left( i\betag(\omega)z-i\omega t +i\phi+ i\phi_0\right) \\ \nonumber
    \E_r'(z,t) &= \real\left\{  \E'(z,t)   \right\}.
\end{align}

此外，由于振荡实电场的瞬时功率与其平方成正比，因此一个周期内的平均功率既可以通过积分求得，也可以直接由复电场绝对值平方的一半得到：
\begin{align}
\label{eq:average_power}
    \langle \E_r^2 \rangle_T&= \frac{1}{T}\int_{0}^{T} \E_r^2 dt = \langle\real\left\{\E\right\}^2\rangle_T =\frac{1}{4}\langle\left(\E+\E^*\right)^2\rangle_T\\ \nonumber
&=\frac{1}{4}\langle \E^2+\E^{*2}+2|\E|^2 \rangle_T=\frac{1}{2} \langle|\E|^2\rangle_T=\frac{1}{2}|\E|^2.
\end{align}
既然复场表示既能避免“手动插入”相位，也能省去显式积分，那么在多数计算中直接使用复电场，并只在必要时取其实部，是非常划算的。因此，本教程主要使用复场表述；但也始终强调，这只是便于计算的抽象，而真正具有物理意义的是实电场，因为它才直接决定电荷的运动。

\subsection{实际电场与电场包络}
Wi-Fi 所用无线信号的载波频率大约是 5~GHz，而先进示波器大约能测到 100~GHz 的电场振荡。相比之下，激光脉冲的电场通常在 100~THz 以上的载波频率振荡。因此，直接用实际复电场 $\E(z,t)=\E_0 \exp\left( i\betag(\omega_0)z-i\omega_0 t\right)$ 来做计算或展示结果往往并不方便，因为这种每秒振荡 $\omega_0/2\pi$ 次的快速变化本来就无法被直接探测。于是，我们通常定义复电场的包络为
\begin{align}
\label{eq:envelope}
    \A(z,t) &=a\cdot\E(z,t)\cdot e^{-i(\betag(\omega_0)z-\omega_0t)},
\end{align}
并以它作为计算对象。其中 $a=\sqrt{0.5\epsilon_0ncA_{eff}}$，$\epsilon_0$ 为真空介电常数，$n$ 为介质折射率，$c$ 为光速，$A_{eff}$ 为光场横截面的有效面积，它们共同给出了归一化因子。用 $a$ 去缩放 $\E$ 可以保证 $\A$ 的量纲是 $\sqrt{W}$，从而 $|\A|^2$ 的量纲就是 $W$。通过把快速、不可直接探测、但又\emph{可预知}的时间和空间振荡“提取”出去，分析电场因线性和非线性效应而发生的\emph{变化}就会容易得多。关于 $\E$ 与 $\A$ 差别的示意，可见 \href{https://www.desmos.com/calculator/rsw2fn5af6 }{这个交互图}。



\section{傅里叶变换}
%Define the Fourier Transform so going from time to frequency and back is well-behaved.

在本教程中，傅里叶变换及其逆变换定义为
\begin{align}
    \Tilde{\E}(z,\omega) &= \FT\left\{\E(z,t)\right\} = \int_{-\infty}^{\infty} \E(z,t) e^{i\omega t} dt, \\ \nonumber
    \E(z,t) &= \IFT\left\{\Tilde{\E}(z,\omega)\right\} = \frac{1}{2\pi} \int_{-\infty}^{\infty} \Tilde{\E}(z,\omega) e^{-i\omega t} d\omega.
\end{align}
这里选择在傅里叶变换中使用 $\exp(i\omega t)$，而不是 $\exp(-i\omega t)$，是因为沿正 $z$ 方向传播的复平面波写作 $\exp(i\betag(\omega_0)z-i\omega_0 t)$。于是有

\begin{align}
    \FT\left\{\exp(i\betag(\omega_0)z-i\omega_0 t)\right\} &= \int_{-\infty}^{\infty} e^{i\betag(\omega_0)z-i\omega_0 t} e^{i\omega t} dt, \\ \nonumber
      &= \int_{-\infty}^{\infty} e^{i\betag(\omega_0)z-i(\omega_0-\omega) t} dt \\ \nonumber
      &= e^{i\betag(\omega_0)z}\delta(\omega-\omega_0),
\end{align}
这说明复平面波的傅里叶变换会在正的载波角频率 $\omega_0$ 处产生一个 delta 函数。如果反过来采用 $\exp(-i\omega t)$ 的定义，那么对沿正 $z$ 方向传播的复平面波做变换时就会出现 $\delta(\omega_0+\omega)$，意味着 delta 函数会落在负载波频率上。对于涉及沿 $z$ 方向传播复平面波的推导，这种写法更容易在符号上出错，因此这里采用前一种约定。

\section{脉冲}
在介质中持续时间有限的电磁脉冲，可以看成无穷多个不同复平面波的叠加：
\begin{align}
    \label{eq:pulse}
    \E(z,t) &= \frac{1}{2\pi}\int_{-\infty}^{\infty} |\Tilde{\E}(z,\omega)| e^{i\betag(z,\omega)z-i\omega t+i\phi(z,\omega)} d\omega \\ \nonumber
    \E(z,t) &= \frac{1}{2\pi}\int_{-\infty}^{\infty} \Tilde{\E}(z,\omega) e^{i\betag(z,\omega)z-i\omega t} d\omega.
\end{align}
式~\ref{eq:pulse} 的关键启发在于：一个光信号在强度、形状和颜色上的任何变化，都必然来自于改变构成它的时间频率分量的振幅 $|\Tilde{\E}(z,\omega)|$、相位 $\phi(z,\omega)$，或者空间频率 $\betag(z,\omega)$。即便是那些看起来很“神奇”的非线性效应，本质上也不过是在改动这三样东西。



\section{啁啾与延迟}
%Highlight that how the phase changes with time determines instantaneous frequency, which is super important for understanding pulse evolution. Explain that the change in phase w.r.t. frequency determines the time shift for each frequency.
考虑在 $z=0$ 处的复电场
\begin{align}
\label{eq:chirp_example}
    a\E(t) &= \A(t)\exp\left(  -i\left(\omega_0 +\frac{C}{2}T \right)T   \right).
\end{align}
如果 $C=0$，那么电场相位随时间线性变化，对应固定的载波频率 $\omega_0$。若 $C>0$，式~\ref{eq:chirp_example} 表明载波频率会随时间升高；若 $C<0$，则说明载波频率会随时间降低。可以参考 \href{https://www.desmos.com/calculator/gd7s8nhfdn}{这个交互图}。电场的“瞬时角频率”定义为
\begin{align}
\label{eq:chirp_definition}
    \delta\omega(T) &= -\partial_T\phi(T),
\end{align}
其中 $\phi(T)$ 是场的相位随时间的函数。时间导数前的负号是为了保证对沿 $z$ 方向传播、写作 $\exp(i\betag(\omega_0)z-i\omega_0 t)$ 的复平面波计算瞬时角频率时，结果能正确得到 $+\omega_0$。

若一个电场的瞬时频率会随时间变化，就称该场带有“啁啾”。当某一时刻的频率低于（高于）载波频率时，可称该处为“红啁啾”（“蓝啁啾”）。如果啁啾从“红”逐渐变到“蓝”，就称其为“递增啁啾”；反之，从“蓝”变到“红”则称为“递减啁啾”。

理解“同一个已知电场的不同时间位置可以具有不同瞬时频率”，并且知道这些频率可以通过式~\ref{eq:chirp_definition} 中对相位取负时间导数得到，对于理解许多线性和非线性效应都非常关键。举例来说，如果某种介质中的光速使得高频光（也就是“更蓝”的光）比低频光（“更红”的光）传播更快，那么一个脉冲在该介质中传播时，其前沿会形成蓝啁啾，而后沿会形成红啁啾。

类似地，时间相位对时间求导给出瞬时频率，频谱相位对频率求导则能告诉我们不同频率分量的时间延迟。考虑在 $z=0$ 处复包络的频谱
\begin{align}
\label{eq:spectrum_time_example}
    \Tilde{\A}(\omega) &= \Tilde{\A}_0(\omega)\exp\left( i\left(\frac{B}{2}(\omega-\omega_0)^2 \right)   \right).
\end{align}
假设 $B>0$，式~\ref{eq:spectrum_time_example} 表明围绕载波角频率 $\omega_0$ 的频率分量，其相位会随着偏离载波的程度二次增长。换个角度说：对于低于 $\omega_0$ 的角频率，相位随频率增加而\emph{下降}；而对于高于 $\omega_0$ 的角频率，相位则随频率增加而\emph{上升}。受式~\ref{eq:chirp_definition} 启发，我们定义

\begin{align}
\label{eq:delay_definition}
    \delta t(\omega) &= \partial_\omega\phi(\omega),
\end{align}
代入式~\ref{eq:spectrum_time_example} 可得
\begin{align}
    \delta t(\omega) &=  B(\omega-\omega_0),
\end{align}
这表明低于 $\omega_0$ 的频率分量会获得负的时间位移（即更早到达），而高于 $\omega_0$ 的频率分量会获得正的时间位移（也就是延迟）。需要注意的是，在式~\ref{eq:delay_definition} 中无需再额外改符号即可保持一致性。

与式~\ref{eq:chirp_definition} 相比，式~\ref{eq:delay_definition} 在实际计算中不那么常用；但“相位随角频率减小意味着提前到达，相位随角频率增大意味着延后到达”这一认识，在第~\ref{ch:Dispersion} 章分析色散时非常有帮助。关于谱相位变化与时间延迟关系的图示，可参考 \href{https://youtu.be/E3S0BQiy3p8}{这个视频教程}。
